\documentclass[11pt]{article}

\usepackage[margin=1in]{geometry}
\usepackage{amsmath,amssymb}
\usepackage{microtype}
\usepackage[hidelinks]{hyperref}

\newcommand{\E}{\mathbb{E}}
\newcommand{\littleo}{o}

\title{The Missing Invariant in the Uniform-Equilibrium Program\\
\large From Incentive Debt to an Escaping Middle}
\author{}
\date{}

\begin{document}

\maketitle

\begin{abstract}
The general uniform-equilibrium conjecture is still open.  Its verification
side is well developed: once a strategically credible public architecture is
supplied, continuation potentials and sublinear error bounds can certify it.
The missing part is a producer that makes local incentive calculations
cooperate along one history-dependent play, including histories generated by
an arbitrary unilateral deviator.  In its shortest form, the desired
principle is:
\begin{quote}
Every incentive debt must either be absorbed by continuation value, paid for
by an observable charge whose cumulative cost is sublinear, or accompanied
by a strict decrease in a finite complexity rank.
\end{quote}
This note explains that sentence without relying on the surrounding analytic
and algebraic machinery.  It then describes the sharper finite-quitting-game
frontier.  There the proof has reached an ``escaping middle'': finite blocks
are understood exactly, but their lengths diverge.  Absorption mass, rather
than calendar time, is the leading proposed coordinate for the limit.  The
required enriched compactification and its strategic decoders remain open.
\end{abstract}

\section{What uniformity asks for}

Consider a finite stochastic game and a proposed long-run payoff vector
\(v\).  For every accuracy \(\varepsilon>0\), a uniform equilibrium payoff
requires a single behavior profile \(\sigma\) and a single threshold
\(T_0\) such that, for every horizon \(T\geq T_0\),
\begin{enumerate}
  \item the average payoff under \(\sigma\) is within \(\varepsilon\) of
        \(v\); and
  \item no player can gain more than \(\varepsilon\) by replacing their
        entire behavior strategy.
\end{enumerate}
The profile may depend on \(\varepsilon\), but it may not depend on the
eventual horizon \(T\).  It must also survive deviations that depend on the
complete observed history.

It is useful to distinguish three evidentiary levels throughout this note.
Statements described as \emph{machine-checked} are present in the production
Lean development.  Statements attributed to the absorption-path literature
are published results under the definitions used there; importing them into
the repository's terminal all-behavior semantics still requires an explicit
adapter.  Objects described as \emph{proposed}, \emph{candidate}, or
\emph{required} are open research obligations, not consequences of
compactness already proved.

Discounted games provide excellent local approximations.  For every discount
rate \(\lambda\), one can choose a stationary discounted equilibrium and its
continuation values.  When \(\lambda\) is small, the players are very patient,
so this looks close to a long-run solution.  But one cannot simply use
\(\lambda=1/T\): the strategy is fixed before the evaluation horizon is
revealed.

The natural response is to let the effective discount index, or more
generally the \emph{public phase}, change with the realized history.  A phase
may record the current state, an account balance, a punishment mode,
transition evidence, or the current level of a finite obstruction
hierarchy.  Each phase can be locally well justified.  The difficulty is
that moving among locally justified phases need not be globally justified.

\section{The no-unpaid-debt ledger}

The central mental model is a ledger.  After every public history, associate
to each player \(i\):
\begin{itemize}
  \item a current public phase;
  \item a continuation credit or potential \(\Phi_i\);
  \item a cumulative error budget \(B_i\); and
  \item a finite rank \(R\) measuring unresolved structural complexity.
\end{itemize}

For obedient play, the desired finite-horizon estimate has the schematic
form
\begin{equation}
  \E\!\left[\sum_{t<T} g_i(t)\right]
  \geq
  T v_i+
 \E[\Phi_i(T)]-\Phi_i(0)-B_i(T).
  \label{eq:lower-ledger}
\end{equation}
For an arbitrary unilateral deviation by player \(i\), one needs the reverse
kind of estimate:
\begin{equation}
  \E\!\left[\sum_{t<T} g_i^{\mathrm{dev}}(t)\right]
  \leq
  T v_i+
 \text{controlled continuation terms}
  +B_i^{\mathrm{dev}}(T).
  \label{eq:upper-ledger}
\end{equation}

The continuation terms telescope: credit granted at one step is accounted
for at the next.  If the surviving boundary terms are controlled and
\[
  B_i(T)=\littleo(T),
  \qquad
  B_i^{\mathrm{dev}}(T)=\littleo(T),
\]
then division by \(T\) makes all ledger errors vanish.  Equations
\eqref{eq:lower-ledger} and \eqref{eq:upper-ledger} then give the uniform
Ces\`aro payoff and deviation bounds.

This verification principle is already formalized.  The missing mathematics
is a construction of phases, potentials, and charges for which the
sublinear-budget claims remain valid under every deviation-induced history
law.

\section{The local trichotomy that must become global}

The asymptotic discounted-equilibrium machinery examines what happens as
patience grows.  Locally, it produces a response of one of the following
forms.

\begin{enumerate}
  \item \textbf{Potential resolution.}
        A harmonic or continuation adjustment absorbs the apparent
        discrepancy.  Nothing observable needs to happen.

  \item \textbf{Observable charge.}
        The obstruction exposes a payoff advantage, an occupation
        discrepancy, or a difference between the baseline and deviating
        transition laws.  This yields a bounded public statistic with zero
        baseline expectation and positive expectation under the selected
        deviation.

  \item \textbf{Finite-rank descent.}
        The obstruction permanently simplifies the remaining problem: an
        action support shrinks, a transient state is removed, a harmonic
        dimension is resolved, or a lexicographic level is completed.
\end{enumerate}

The desired invariant is therefore the global alternative
\[
  \boxed{
    \text{resolved by potential}
    \quad\lor\quad
    \text{paid by sublinear charge}
    \quad\lor\quad
    \text{strict finite-rank descent}.
  }
\]
No local incentive discrepancy may disappear merely because the strategy
changed its phase or selected a new discounted equilibrium.

Rank descent can occur only finitely often.  Charges may occur repeatedly,
but their expected cumulative total must be sublinear in the horizon.
Together, these conditions rule out endless cost-free refinement.

\section{Why a pointwise detector is not an invariant}

Suppose a deviation changes the next-state distribution at every round.
Because the state space is finite, at each round some signed state indicator
detects the difference.  It is tempting to run all such indicators through
an online learner and conclude that the deviation will eventually be
detected.

The conclusion is false without an additional stability principle.  The
revealing coordinate may alternate.  On odd rounds, destination \(A\) may
give the positive score; on even rounds, destination \(B\) may give it.  A
positive detector then exists at every round even though the cumulative
score of every fixed detector is zero.  The local positive signals cancel.

There is an even stronger difficulty.  Once the public strategy announces a
mixture of monitors, an adaptive deviator may be able to choose a different
transition law whose expected score against that mixture is nonpositive.
Thus neither a finite-family pigeonhole argument nor ordinary fixed-action
regret supplies the missing global conclusion.

A valid invariant must instead ensure at least one of the following:
\begin{itemize}
  \item one analytic branch and one detector remain fixed on a sufficiently
        stable block;
  \item changes of detector have a quantitatively bounded switching cost;
  \item every change is charged to observable deviation evidence; or
  \item every change strictly decreases the finite rank.
\end{itemize}
The important object is not ``a detector exists now,'' but ``the detector can
accumulate evidence before it is changed, and changing it is never free.''

\section{Why occupation control is not enough}

A second plausible invariant tracks how much time the process spends in
transient states.  This is also insufficient.

A deviator may use one transient move to redirect play into a different
recurrent class---for example, into a bad absorbing state rather than a good
one.  The transient discrepancy then lasts only one round, while the payoff
difference after entry is linear in \(T\).  A small transient occupation
budget therefore does not imply a small payoff error.

The invariant must remember more than the amount of transient occupation.  It
must preserve:
\begin{itemize}
  \item which recurrent component is entered;
  \item which obeying and deviating ``shadow plays'' are being compared; and
  \item which continuation promise is attached to the entry state.
\end{itemize}
In short, it must control continuation identity, not merely occupation mass.

\section{The lesson of the Big Match}

The Big Match shows why ``make the selector history-dependent'' is not a
complete answer.  A natural linear running-surplus selector produces stopping
hazards of harmonic order, roughly \(1/t\).  Their sum diverges, so stopping
eventually occurs with overwhelming probability in precisely the histories
where continued survival is needed.  The resulting average payoff can be
wrong even though every individual discounted choice looks reasonable.

The successful Blackwell--Ferguson mechanism uses a square-order hazard,
roughly \(1/D^2\), tied to a nonlinear account.  Its occupation and
summability balance is different.  The example demonstrates that a global
invariant must control not only which phase is selected, but also the total
mass spent on rare transitions, resets, and phase changes.

\section{Why the multiplayer case is coupled}

In a two-player zero-sum game, a strategy securing one player's value has a
direct relationship to the other player's upper bound.  This permits an
account construction in which continuation credit, value switching, and
floor occupation can be analyzed through one-sided guarantees.

With several players, punishment promises interact.  A continuation that
discourages player \(i\) may make deviation attractive to player \(j\), who
is supposed to carry out the punishment.  Separately feasible playerwise
continuation inequalities need not combine into one feasible public system.

Consequently, the missing invariant must be simultaneous:
\begin{itemize}
  \item cooperative and punishment phases use one public update rule;
  \item every player's unilateral deviation is capped;
  \item the players implementing a punishment are themselves willing to
        follow it; and
  \item any incompatibility among continuation requirements creates a new
        charged response or a strict rank decrease.
\end{itemize}
This coupled requirement contains much of the genuinely multiplayer
difficulty.

\section{The sharper quitting-game front}

Finite quitting games are presently the most direct laboratory for this
problem.  There is one live state.  At each visit, every player chooses
whether to continue or quit; a nonempty set of quitters ends the game and
selects a terminal payoff.  This removes recurrent-state geometry without
removing the multiplayer credibility problem.

Several reductions on this front are machine-checked.  Existence of terminal
approximate Nash profiles at every accuracy is equivalent to existence of a
uniform-equilibrium payoff.  Against fixed opponents, an arbitrary behavioral
deviation is payoff-equivalent to a mixture of deterministic quit times and
the option of never quitting.  Finite Nash--Bellman chains can be optimized at
each cutoff, and their optimized dynamic debt has an exhaustive zero-versus-
positive limit split.  The zero-debt branch already produces terminal
approximate equilibria.  Thus the live branch is a positive debt plateau, not
an unspecified failure of compactness.  This split is exhaustive for the
selected exact-chain optimization; it is not a claim that every possible
equilibrium belongs to that grammar.

The plateau carries much more information than one positive number.  It
selects a debt owner, summable opponent-only survival clocks, and a marked
terminal quitter set at the reverse end.  Every finite middle block has an
exact associative summary: prescribed payoff transport, unilateral stopping
data, its entry and exit roots, and the provenance of the marked packet.
Fixed-cutoff compactness of this full finite data is also machine-checked.
What escapes is the number of stages in the middle.

This is an exhaustive reduction inside finite quitting games; it is not a
reduction of arbitrary finite stochastic games to quitting games.  A positive
solution would settle the quitting subclass.  A quitting counterexample would
refute the general universal conjecture, but solving quitting games would not
by itself prove the general conjecture.

\section{Why calendar time cannot carry the limit}

Suppose a block has length \(m\), and attach to it all of its bounded payoff,
survival, obstacle, and holonomy coordinates.  Keeping \(m\) literally places
the block in a space of the form
\[
  \mathbb{N}\times X,
\]
where \(X\) is compact.  Every compact subset of this product has bounded
\(\mathbb{N}\)-coordinate.  Consequently, a family whose middle lengths tend
to infinity cannot have a compact limit while literal game-stage length
remains part of the state.  This elementary obstruction, including its
application to the calibrated quitting blocks, is machine-checked.

Adding a formal point called \(\infty\) repairs the topological statement but
not the game-theoretic one.  A point at infinity is not yet a finite strategy.
To use it, one still needs a theorem decoding it either into an
accuracy-dependent finite behavior profile or into a uniformly bounded
finite modification.  Compact coefficients alone do not pay for the stages
needed to implement them.

The obstruction is especially visible in a low-hazard interval.  Let \(q_t\)
be the probability of absorption at stage \(t\), conditional on surviving to
that stage, and let
\[
  s_t=\prod_{u<t}(1-q_u),
  \qquad
  a_t=s_tq_t.
\]
The calendar interval may contain more and more stages while each \(q_t\)
shrinks.  Its strategic effect is carried by the absorption masses \(a_t\),
not by the number of nearly inert calendar ticks.  This is why calendar time
is the wrong coordinate for the escaping middle.

\section{Why accumulated absorption mass is the leading coordinate}

An absorption path uses cumulative absorbed mass as its clock.  Long runs of
tiny hazards are compressed, but where absorption occurs and which terminal
payoff it selects remain visible.  Published absorption-path theory proves
compactness and equilibrium connections for its base object under its own
definitions.  It is therefore the leading existing model for the boundary we
actually see, rather than an analogy chosen after the fact.

The repository needs a stricter object, provisionally a \emph{marked
subprobability absorption path}.  It is only a proposed compactification.  In
particular, adding the strategic coordinates below may destroy the closedness
enjoyed by the unmarked object.  A published compactness theorem for ordinary
absorption paths does not prove compactness of this enrichment.

There is a basic distinction already at the level of total mass.  At the end
of a finite block, residual survival probability is an \emph{exit port}.  It
must be multiplied into the next block, where it may later absorb.  At the end
of a completed infinite play, residual survival probability is the genuine
\emph{Never} alternative.  Concatenation relates the two, but they are not the
same atom.  Declaring every finite-block remainder to be Never would make
block composition give the wrong payoff and the wrong deviation values.

\section{What the enriched path must remember}

An ordinary terminal distribution is too coarse for the plateau consumer.
At least four additional kinds of information appear necessary.

\paragraph{Anchors and chronology.}
The path must retain the entry root, the exit continuation, and enough source
data to prove that adjacent pieces can actually be spliced in chronological
order.  A limit of scalar summaries need not be a limit of executable blocks.

\paragraph{The conditional marked packet.}
The positive-debt extraction identifies a complete simultaneous quitter set
and its local incentive comparison at a selected reverse root.  Its mass as
seen from the distant forward entry is multiplied by survival through the
whole middle.  That transported mass can become uninformative even while the
conditional terminal kernel and its incentive direction remain decisive.
The limit therefore needs a pointed or blown-up conditional mark, not merely
the unconditional terminal measure and not a singleton chosen from the
quitter set.

\paragraph{The full stopping obstacle.}
For each player, quitting at each possible live date and never quitting are
the extreme deviation choices.  The relevant datum is therefore the graph of
quit-time payoffs, together with its Snell envelope, rather than only the
largest scalar cap.  A supremum records how much a player can gain but may
forget where the gain is nearly attained, which continuation made the quit
legal, and whether that witness persists under a limiting seam.  Those are
exactly the facts a decoder needs.

\paragraph{Debt and playerwise clocks.}
The enriched path must still identify the debt owner, the owner's and
opponents' survival clocks, the promised payoff path, and the debt measured at
the original entry.  A local improvement at a late root is useless unless its
benefit transports back to that original optimization problem.

The first responsible formalization target is finite, not topological: encode
every calibrated finite block as such a marked absorption cylinder and prove
exact payoff, obstacle, debt, packet, anchor, and concatenation identities.
Only after this semantic map is fixed is it meaningful to choose an infinite
topology.  A decisive falsification experiment would construct two finite
families with the same ordinary absorption-path limit but incompatible
limiting obstacles, conditional packets, or splice behavior.

\section{The two decoder problems}

Even a compact closed enriched path space would not finish the argument.
There are two different consumers, and neither is presently proved.

The first is a \emph{valid-path compiler}.  From a globally admissible limit
path it must construct, for every accuracy, an actual behavior profile
satisfying terminal payoff and every all-behavior unilateral
deviation inequality.  A sure terminal jump needs special care: the published
sequential-perfection condition does not automatically supply the repository's
required endpoint test.  The adapter must either augment that endpoint or stop
before it and discharge the final continuation by an already verified direct
repair.

The second is a \emph{failed-path surgery}.  If a proposed limit is not
globally admissible, a strict local failure must yield an executable finite
extension of bounded length.  The extension must preserve the original
anchor and marked packet, restore exact Nash--Bellman roots, and decrease the
optimized debt by a fixed amount measured back at the original entry.  Merely
finding a different continuation value, a positive but vanishing improvement,
or a seam in a relaxed coefficient space is not enough.

There is a clean way these local obligations could become uniform.  Fix an
accuracy and a positive plateau level.  If every point \(z\) of a compact
enriched state space has either a stable repair certificate or a neighborhood
on which some finite surgery of length \(L_z\) lowers root debt by
\(c_z>0\), then a finite subcover supplies uniform finite bounds and a uniform
positive decrement.  Equivalently, failure of uniformity produces a
convergent bad sequence contradicted by the local alternative at its limit.
This finite-cover argument is sound only after compactness, closed strategic
semantics, and stability of the local decoders are proved.  It does not evade
the escaping-length obstruction.

The remaining hard work is consequently concrete:
\begin{enumerate}
  \item define and verify the exact finite-cylinder encoding;
  \item prove the marked payoff--obstacle--packet graph sequentially compact
        and closed, or find the missing state by a counterexample;
  \item prove the corrected endpoint adapter and valid-path compiler;
  \item prove a stable bounded surgery with cutoff-independent descent; and
  \item feed either output through the machine-checked terminal-to-uniform
        bridge.
\end{enumerate}
The existing finite-block holonomy, depth-free reinsertion estimate, and
buffered return/exit/dead-end lemmas are inputs to these tasks.  They do not
yet solve either decoder.

\section{What the advanced tools contribute}

The surrounding tools have distinct conceptual jobs.
\begin{description}
  \item[Analytic and Puiseux germs]
    retain the rates at which discounted equilibria move as patience grows.
    Plain convergence would forget exactly the rare-action information
    exposed by Big-Match-type examples.

  \item[Farkas duality]
    turns failure of a desired continuation adjustment into a concrete
    obstruction explaining that failure.

  \item[Transition monitors]
    turn differences between probability laws into bounded public evidence.

  \item[Markov potentials]
    turn occupation and transition discrepancies into telescoping accounts.

  \item[Online learning]
    watches a finite collection of possible monitors without knowing the
    deviator's plan in advance.

  \item[Well-founded ranks]
    ensure that structural refinement cannot continue forever.

  \item[The public-phase compiler]
    converts completed historywise potential and charge inequalities into
    the final uniform-equilibrium bounds.
\end{description}

In their proved scopes, these tools provide components for local extraction,
statistical monitoring, telescoping, and verification.  They do not yet form
one producer.  The missing bridge is the game-specific global rule that
prevents their switching, synchronization, punishment, and implementation
costs from becoming linear.

\section{What changed at the terminal layer}

Since the escaping-middle picture above was drawn, the quitting front has
gained a second, stationary attack, and it is machine-checked where stated.

The price of punishment is now exact.  The worst that opponents can inflict
on one player --- over \emph{all} history-dependent plans and all replies ---
equals the best constant row's simple two-branch stopping value: quit into
the row now, or ride its absorption forever.  Nothing time-varying punishes
one designated player better.  This closes a long-standing gap between the
crude ``solo-clipped'' punishment ceiling and the true floor, and an
explicit table shows the gap was real.

The one-stage mechanisms are now a catalogue rather than a search space.
Which coalitions can quit surely in equilibrium, which owner rates are
blocked and by whom, and exactly when a blocked owner can be repaired by
colliding with a sure blocker --- each has a finite list of scalar
conditions, priced against the exact punishment floor, proved in both
directions.  Two structural lessons came out.  First, rarity buys nothing:
making an invalid repair phase rare cannot fix it, because the phase's
inequalities do not see how often it is played.  Second, when the
catalogued mechanisms fail, the missing move is not a cleverer schedule but
a wider stage: letting a third player mix --- support enlargement --- creates
exact equilibria that no rotation of narrower phases reaches.  What remains
at three coordinates is a single finite semialgebraic decision: either the
residue after the catalogue is empty, or a member survives and must be fed
to the cycle engines.

One surprise is worth a sentence.  Punishing one player is stationary;
punishing \emph{everyone at once} with a single plan is not free, and on an
explicit cyclic table every shared plan leaves someone three quarters above
their floor.  On a sibling table, staggered time-varying plans beat every
constant row through exactly the mechanism of non-transitive dice.  Whether
stationarity genuinely fails for simultaneous punishment is now a posed
question; either answer prices what a single shared punishment tail can
promise.

\section{The invariant in one sentence}

\begin{quote}
\emph{Along every possible deviating history, every local incentive gain
must either be neutralized by continuation value, leave accumulating public
evidence, or permanently simplify the remaining problem; and the total cost
of everything that does not permanently simplify the problem must be
sublinear in time.}
\end{quote}

This is not merely a missing Lean lemma.  In the unrestricted multiplayer
case it is one formulation of the core producer missing from the open
uniform-equilibrium problem.
The existing formal development has isolated the interfaces on both sides:
local discounted and obstruction data enter from one side, and a checked
public-phase verifier waits on the other.  The missing invariant is the
history-safe conservation law that joins them.

\end{document}
